\documentclass{amsart}

\usepackage{amssymb}
\usepackage{mathtools}
\usepackage{booktabs}
\usepackage{enumitem}
\usepackage[hidelinks]{hyperref}

\hypersetup{
  pdftitle={Hilbert coefficients of complete bipartite quadratic ideals},
  pdfauthor={Knowledge Sense},
  pdfsubject={Commutative algebra},
  pdfkeywords={Hilbert coefficient, Hilbert-Samuel polynomial, quadratic algebra, monomial ideal, fiber product, complete bipartite graph}
}

\newtheorem{theorem}{Theorem}[section]
\newtheorem{proposition}[theorem]{Proposition}
\newtheorem{lemma}[theorem]{Lemma}
\newtheorem{corollary}[theorem]{Corollary}
\theoremstyle{definition}
\newtheorem{definition}[theorem]{Definition}
\newtheorem{notation}[theorem]{Notation}
\theoremstyle{remark}
\newtheorem{remark}[theorem]{Remark}
\newtheorem{example}[theorem]{Example}

\numberwithin{equation}{section}

\newcommand{\NN}{\mathbb{N}}
\newcommand{\QQ}{\mathbb{Q}}
\newcommand{\mm}{\mathfrak{m}}
\newcommand{\cf}[1]{[#1]\,}
\newcommand{\fl}[1]{\lfloor #1\rfloor}

\begin{document}

\title[Complete bipartite quadratic ideals]{Hilbert coefficients of complete bipartite quadratic ideals}

\author{Knowledge Sense}
\address{Knowledge Sense}
\email{math@knowledgesense.jp}

\date{12 October 2026}

\subjclass[2020]{Primary 13D40; Secondary 13F55, 05A15, 05E40}

\keywords{Hilbert coefficient, Hilbert--Samuel polynomial, quadratic algebra, monomial ideal, fiber product, complete bipartite graph}

\begin{abstract}
Let $k$ be a field and let $1\le p\le q$. In the polynomial ring $S$ over $k$
in variables $x_1,\dots,x_p,y_1,\dots,y_q$ consider the ideal $I_{p,q}$
generated by the squares of all variables and by all products $x_ay_b$, that
is, the edge ideal of the complete bipartite graph $K_{p,q}$ together with all
squares. We compute the lengths of $S/I_{p,q}^s$ for all $s\ge1$ and all
Hilbert coefficients of $I_{p,q}$. The Hilbert--Samuel polynomial gives
$\ell(S/I_{p,q}^s)$ for every $s\ge1$, the last coefficient $e_{p+q}$
vanishes, and $e_0,\dots,e_p$ coincide with the Hilbert coefficients of the
square of the maximal ideal. Each of the remaining coefficients
$e_{p+1},\dots,e_{p+q-1}$ is nonzero and has sign $(-1)^{p+1}$. Consequently
all Hilbert coefficients are nonnegative if and only if the smaller part size
$p$ is odd. For $p=2$ we give closed formulas, which exhibit $q-1$ consecutive
negative coefficients. The four-cycle $p=q=2$, the smallest member of the
family with a negative coefficient, has $e_3=-2$ and is a counterexample to a
conjecture of Fr\"oberg on the nonnegativity of Hilbert coefficients of
Artinian quadratic algebras. In Fr\"oberg's list of predicted coefficients in
four variables this ideal appears with $e_3=2$; we show that the sign is
negative. The ideal is a fiber product, and the proof starts from the known
decomposition of powers of fiber products.
\end{abstract}

\maketitle

\section{Introduction}

Let $S=k[z_1,\dots,z_n]$ be a polynomial ring over a field $k$, and let
$I\subseteq S$ be a homogeneous ideal such that $S/I$ is Artinian. By a theorem
of Samuel~\cite{Sam} the length $\ell(S/I^s)$ is a polynomial in $s$ of degree $n$ for
large $s$. Following Fr\"oberg~\cite{Fr} we write it as
\begin{equation}\label{eq:convention}
  \ell(S/I^s)=\sum_{i=0}^{n}(-1)^i e_i(I)\binom{s+n-i-1}{n-i}
  \qquad (s\gg 0),
\end{equation}
and call $e_0(I),\dots,e_n(I)$ the \emph{Hilbert coefficients} of $I$. The
same convention defines the Hilbert coefficients of an $\mm$-primary ideal of
a Cohen--Macaulay local ring $(R,\mm)$. For such ideals $e_0$ is positive,
Northcott~\cite{Nor} showed that $e_1\ge0$, and Narita~\cite{Nar} showed that
$e_2\ge0$ and gave an example with $e_3<0$; see also the introductions
of~\cite{Fr} and~\cite{MN}.

The paper~\cite{Fr} studies Hilbert coefficients of Artinian quadratic
algebras. On the basis of extensive calculations, Fr\"oberg conjectured in
this setting that negative coefficients do not occur~\cite[Conjecture~1.1]{Fr}:
if $I$ is generated by quadrics and $S/I$ is Artinian, then $e_i(I)\ge0$ for
all $i$. In the same paper he computed or predicted the Hilbert coefficients
of all Artinian quotients by quadratic monomial ideals in at most four
variables~\cite[Theorem~3.1, Conjecture~3.2]{Fr}, and conjectured a formula for
the Hilbert coefficients of the square of the maximal
ideal~\cite[Conjecture~3.1]{Fr}.

In this paper we study the following family. For integers $1\le p\le q$ put
$n=p+q$ and
\[
  S=k[x_1,\dots,x_p,y_1,\dots,y_q],\qquad
  I_{p,q}=\bigl(x_a^2,\ y_b^2,\ x_ay_b \;:\; 1\le a\le p,\ 1\le b\le q\bigr).
\]
Thus $I_{p,q}$ is the edge ideal of the complete bipartite graph $K_{p,q}$
together with the squares of all variables; in the notation of~\cite{Fr} it is
the ideal $I_{n,M}$ with $M$ the edge set of $K_{p,q}$. Let
$\mm_x=(x_1,\dots,x_p)$, $\mm_y=(y_1,\dots,y_q)$,
$Q_x=(x_1^2,\dots,x_p^2)$ and $Q_y=(y_1^2,\dots,y_q^2)$. Then
\[
  I_{p,q}=Q_x+Q_y+\mm_x\mm_y
\]
is the \emph{fiber product} of $Q_x\subseteq k[x]$ and $Q_y\subseteq k[y]$ in the
sense of Nguyen and Vu~\cite{NV} and of Do, Nguyen and Seyed
Fakhari~\cite{DNSF}; correspondingly $S/I_{p,q}\cong A\times_k B$ with
$A=k[x]/Q_x$ and $B=k[y]/Q_y$ \cite[Section~1]{DNSF}.

To state the result we use the following notation. For a power series $f$ and
an integer $r$, $\cf{z^r}f$ denotes the coefficient of $z^r$ in $f$, which is
$0$ if $r<0$. For integers $a,b\ge1$ and $s\ge1$ let
\begin{equation}\label{eq:D}
  D_{a,b}(s)=\sum_{\substack{t\ge0,\ r\ge1\\ t+r\le s}}\ \sum_{\kappa=t+2r}^{a}
  \binom{t+b-1}{b-1}\binom{a}{\kappa}\binom{s-t-r+a-1}{a-1},
\end{equation}
where an empty inner sum is $0$, and for integers $m\ge0$ and $i\ge1$ let
\begin{equation}\label{eq:Delta}
  \Delta_{m,i}(z)=\sum_{j=m}^{i-1}(-1)^j\binom{i-1}{j}(1-z)^j,
\end{equation}
which is $0$ if $m\ge i$.

\begin{theorem}\label{thm:main}
Let $k$ be a field, let $1\le p\le q$, $n=p+q$ and $I=I_{p,q}$.
\begin{enumerate}[label=\textup{(\alph*)}]
\item For every integer $s\ge1$,
\[
  \ell(S/I^s)=\binom{2s+n-1}{n}+D_{p,q}(s)+D_{q,p}(s).
\]
\item Equation~\eqref{eq:convention} holds for every integer $s\ge1$, with
$e_0(I)=2^n$ and
\[
  e_i(I)=\cf{z^{n-i-1}}\frac{z^{i-1}-\Delta_{p,i}(z)-\Delta_{q,i}(z)}{(1-2z)^{i+1}}
  \qquad (1\le i\le n).
\]
In particular $e_i(I)=2^{n-2i}\binom{n-i}{i}>0$ for $0\le i\le p$, and
$e_n(I)=0$.
\item For $p+1\le i\le n-1$ one has $(-1)^{p+1}e_i(I)>0$.
\item All Hilbert coefficients of $I$ are nonnegative if and only if $p$ is odd.
If $p$ is even, then $e_{p+1}(I),\dots,e_{n-1}(I)$ are $q-1\ge1$ consecutive
negative coefficients.
\end{enumerate}
\end{theorem}

For $p=q=1$ the range in~(c) is empty, and $I_{1,1}$ is the square of the
maximal ideal of $k[x_1,y_1]$. For $p=2$ the formula in~(b) can be made
explicit.

\begin{corollary}\label{cor:k2q}
Let $q\ge2$, $n=q+2$ and $I=I_{2,q}$. Then $e_0(I)=2^n$,
$e_1(I)=(n-1)2^{n-2}$, $e_2(I)=\binom{n-2}{2}2^{n-4}$, $e_n(I)=0$,
\[
  e_{n-1}(I)=3-n+(-1)^{n-1},
\]
and for $3\le i\le n-2$, with $r=n-i-1\ge1$,
\begin{align*}
  e_i(I)&=(2-i)\,2^r\binom{r+i}{i}+(i-1)\,2^{r-1}\binom{r+i-1}{i}\\
  &=-\frac{2^{r-1}}{i}\binom{r+i-1}{i-1}\bigl(2i(i-2)+(i-3)r\bigr).
\end{align*}
In particular $e_3(I),\dots,e_{n-1}(I)$ are $n-3$ consecutive negative
coefficients.
\end{corollary}

By Theorem~\ref{thm:main}(d), a negative coefficient requires $p$ even, hence
$p,q\ge2$ and $n\ge4$. The smallest member of the family with a negative
Hilbert coefficient is therefore the four-cycle $C_4=K_{2,2}$.

\begin{corollary}\label{cor:c4}
For $p=q=2$ and every $s\ge1$,
\[
  \ell(S/I_{2,2}^s)=\binom{2s+3}{4}+2s
  =16\binom{s+3}{4}-12\binom{s+2}{3}+\binom{s+1}{2}+2\binom{s}{1},
\]
so that $(e_0,e_1,e_2,e_3,e_4)=(16,12,1,-2,0)$. In particular $I_{2,2}$ is a
counterexample to \cite[Conjecture~1.1]{Fr}.
\end{corollary}

The first lengths are $7$, $39$ and $132$. Written in four variables, the
ideal of Corollary~\ref{cor:c4} is the case
$M=\{x_1x_2,x_2x_3,x_3x_4,x_4x_1\}$ of~\cite[Conjecture~3.2]{Fr}. There Fr\"oberg predicts the length formula
above and lists the Hilbert coefficients as $(16,12,1,2)$. Under the
convention~\eqref{eq:convention}, which is the one of~\cite[p.~315]{Fr}, the
predicted lengths give $e_3=-2$, so the listed value of $e_3$ has the wrong
sign. The predicted lengths themselves are correct by
Theorem~\ref{thm:main}(a).

\subsection*{Relation to earlier work}
The fiber-product structure of $I_{p,q}$ is the starting point of the proof.
The powers of a fiber product $F=I+J+\mm\mathfrak n$ were decomposed by Nguyen
and Vu~\cite[Proposition~4.4]{NV}, and the resulting expansion
$F^s=\sum_{j=0}^{s}(\mm\mathfrak n)^{s-j}(I^j+J^j)$ was recorded by Do, Nguyen
and Seyed Fakhari~\cite[Lemmas~2.10 and~5.6]{DNSF}; we use it in the form of
Lemma~\ref{lem:powers}. What is new here is the explicit evaluation of all
lengths and of all Hilbert coefficients for the family $I_{p,q}$, and the sign
classification.

Within Fr\"oberg's tables, the family $I_{p,q}$ contains the square of the
maximal ideal in two variables ($p=q=1$), the case $M=\{x_1x_2,x_1x_3\}$ of
\cite[Theorem~3.1]{Fr} ($p=1$, $q=2$), and the cases $M=\{x_1x_2,x_1x_3,x_1x_4\}$
and $M=\{x_1x_2,x_2x_3,x_3x_4,x_4x_1\}$ of \cite[Conjecture~3.2]{Fr}
($p=1$, $q=3$ and $p=q=2$); see Remark~\ref{rem:froberg}. The formula for
$e_0,\dots,e_p$ in Theorem~\ref{thm:main}(b) comes from the square of the
maximal ideal. For this ideal Proposition~\ref{prop:square} proves the
coefficient formula of~\cite[Conjecture~3.1]{Fr}, together with the
accompanying binomial identity after a correction of one index; see the
discussion following Proposition~\ref{prop:square}.

Negative higher Hilbert coefficients, also in consecutive positions, are known
for $\mm$-primary ideals of regular local rings; the examples of Mafi and
Naderi~\cite[Examples~2.1 and~2.2]{MN} have cubic generators or generators of
mixed degrees. Theorem~\ref{thm:main} concerns quadratic ideals: for the family
$I_{p,q}$ it determines the signs of all Hilbert coefficients, by the parity of
one part size.

Banerjee, Goel and Verma~\cite[Theorem~3.3, Example~3.11]{BGV} computed lengths
attached to edge ideals of complete bipartite graphs, but for a different
object; see Remark~\ref{rem:bgv}.

\subsection*{Dependencies}
The decomposition of powers in Lemma~\ref{lem:powers} is due to~\cite{NV,DNSF};
we include a short proof under the hypotheses used here. The rest of the proof
is an exact count of standard monomials and a manipulation of formal power
series with rational coefficients.

\subsection*{Organization}
Section~\ref{sec:prelim} fixes conventions and derives a membership criterion
for powers of $I_{p,q}$ from the fiber-product decomposition.
Section~\ref{sec:lengths} counts standard monomials and proves
Theorem~\ref{thm:main}(a) together with the formula for $e_0,\dots,e_p$.
Section~\ref{sec:coefficients} proves the coefficient formula of
Theorem~\ref{thm:main}(b), and Section~\ref{sec:signs} determines the signs.
Section~\ref{sec:examples} treats the cases $p=1$ and $p=2$, compares with the
literature and lists the coefficients for small $p$ and $q$.

\section{Preliminaries}\label{sec:prelim}

\subsection{Lengths and Hilbert coefficients}
All binomial coefficients $\binom{N}{c}$ below have $N\ge0$, and they vanish
unless $0\le c\le N$. For a power series $f$ we write
$\cf{z^r}f$ as in the introduction. The following lemma connects the
convention~\eqref{eq:convention} with generating functions.

\begin{lemma}\label{lem:numerator}
Let $n\ge1$ and let $L\colon\{1,2,\dots\}\to\QQ$ be a function with
\[
  \sum_{s\ge1}L(s)z^s=\frac{z\,h(z)}{(1-z)^{n+1}}
\]
for a polynomial $h\in\QQ[z]$ of degree at most $n$. Write
$h(1+u)=\sum_{i=0}^{n}e_iu^i$. Then
\[
  L(s)=\sum_{i=0}^{n}(-1)^ie_i\binom{s+n-i-1}{n-i}\qquad\text{for every }s\ge1.
\]
\end{lemma}

\begin{proof}
We have $h(z)=\sum_i e_i(z-1)^i=\sum_i(-1)^ie_i(1-z)^i$, hence
$zh(z)/(1-z)^{n+1}=\sum_i(-1)^ie_i\,z/(1-z)^{n-i+1}$, and the coefficient of
$z^s$ in $z/(1-z)^{n-i+1}$ is $\binom{s-1+n-i}{n-i}$ for every $s\ge1$.
\end{proof}

Since the polynomials $\binom{s+n-i-1}{n-i}$, $0\le i\le n$, form a basis of the
polynomials in $s$ of degree at most $n$, the numbers $e_i$ in
Lemma~\ref{lem:numerator} are uniquely determined by $L$. If
$L(s)=\ell(S/I^s)$ for an ideal $I$ as in~\eqref{eq:convention}, they are the
Hilbert coefficients of $I$, and~\eqref{eq:convention} holds for every $s\ge1$.

For a monomial ideal $J\subseteq S$ with $S/J$ Artinian, $\ell(S/J)$ is the
number of monomials not in $J$, the \emph{standard monomials} of $J$. Powers of
monomial ideals are monomial ideals, and membership of a monomial in a monomial
ideal is a combinatorial condition; hence the lengths in Theorem~\ref{thm:main}
do not depend on $k$. They also do not change if $S$ is replaced by its localization
at the homogeneous maximal ideal.

\subsection{Monomials}
From now on $1\le p\le q$, $n=p+q$, $S=k[x_1,\dots,x_p,y_1,\dots,y_q]$ and
$I=I_{p,q}$. We write $\NN=\{0,1,2,\dots\}$. Every monomial of $S$ is $x^\alpha y^\gamma$ with $\alpha\in\NN^p$
and $\gamma\in\NN^q$. For $\alpha\in\NN^p$ we put
\[
  |\alpha|=\sum_a\alpha_a,\qquad
  \omega(\alpha)=\#\{a:\alpha_a\text{ odd}\},\qquad
  \fl{\alpha/2}=\sum_a\fl{\alpha_a/2}=\frac{|\alpha|-\omega(\alpha)}{2},
\]
and similarly for $\gamma\in\NN^q$.

\subsection{Powers of fiber products}
Let $R$ and $T$ be polynomial rings over $k$ of positive dimension, with
homogeneous maximal ideals $\mm$ and $\mathfrak n$, and let $I'\subseteq\mm^2$
and $J'\subseteq\mathfrak n^2$ be homogeneous ideals. We identify ideals of $R$
and of $T$ with their extensions to $R\otimes_kT$, and we use the convention
that the zeroth power of an ideal is the unit ideal. The fiber product of $I'$
and $J'$ is $F=I'+J'+\mm\mathfrak n$~\cite[Notation~4.3]{NV}.

\begin{lemma}[Nguyen--Vu, Do--Nguyen--Seyed Fakhari]\label{lem:powers}
In this situation, for every $s\ge1$,
\[
  F^s=\sum_{j=0}^{s}(\mm\mathfrak n)^{s-j}\bigl(I'^{\,j}+J'^{\,j}\bigr).
\]
\end{lemma}

\begin{proof}
Put $P=\mm\mathfrak n$. Since $F=I'+J'+P$, the ideal $F^s$ is the sum of the
ideals $I'^{\,a}J'^{\,b}P^c$ with $a+b+c=s$. If $a,b\ge1$, then
$I'J'\subseteq\mm^2\mathfrak n^2=P^2$ gives
$I'^{\,a}J'^{\,b}P^c\subseteq I'^{\,a-1}J'^{\,b-1}P^{c+2}$, and the exponents
on the right again sum to $s$. Repeating this step, every summand is contained
in one with $a=0$ or $b=0$, that is, in $P^{s-j}I'^{\,j}$ or $P^{s-j}J'^{\,j}$
for some $0\le j\le s$. These ideals are themselves summands of $F^s$.
\end{proof}

Lemma~\ref{lem:powers} is \cite[Lemma~2.10]{DNSF}, and it also follows from
\cite[Proposition~4.4(1)]{NV}, namely $F^s=H^s+\sum_{i=1}^{s}(\mm\mathfrak
n)^{s-i}J'^{\,i}$ with $H=I'+\mm\mathfrak n$, combined with the binomial
expansion of $H^s$. In arXiv:2409.17551v1, \cite[Lemma~2.10]{DNSF} is stated
for arbitrary proper homogeneous ideals, but its proof uses the inclusion
$I'J'\subseteq\mm^2\mathfrak n^2$, which holds under the hypotheses above. We
use the lemma only under these hypotheses. Some such hypothesis is needed: for
$R=k[x]$, $T=k[y]$, $I'=(x)$ and $J'=(y)$ one has $F=(x,y)$, and $xy\in F^2$
does not lie in the right-hand side for $s=2$.

We apply Lemma~\ref{lem:powers} with $R=k[x]$, $T=k[y]$, $I'=Q_x$ and $J'=Q_y$,
so that $F=I_{p,q}$.

\begin{lemma}\label{lem:membership}
Let $s\ge0$. A monomial $x^\alpha y^\gamma$ lies in $I_{p,q}^s$ if and only if
\[
  |\alpha|+|\gamma|\ge2s,\qquad
  \fl{\alpha/2}+|\gamma|\ge s,\qquad
  \fl{\gamma/2}+|\alpha|\ge s.
\]
\end{lemma}

\begin{proof}
For $s=0$ there is nothing to prove, so let $s\ge1$. By
Lemma~\ref{lem:powers}, $I^s$ is the sum of the monomial ideals
$\mm_x^{s-j}Q_x^{\,j}\,\mm_y^{s-j}$ and $\mm_x^{s-j}\,\mm_y^{s-j}Q_y^{\,j}$,
$0\le j\le s$. A monomial lies in a sum of monomial ideals if and only if it lies
in one of them. If $\mathfrak a\subseteq k[x]$ and $\mathfrak b\subseteq k[y]$
are monomial ideals, then $x^\alpha y^\gamma\in\mathfrak a\mathfrak b$ if and
only if $x^\alpha\in\mathfrak a$ and $y^\gamma\in\mathfrak b$, because
$\mathfrak a\mathfrak b$ is generated by products of a generator of $\mathfrak
a$ and a generator of $\mathfrak b$. Moreover $y^\gamma\in\mm_y^{s-j}$ if and
only if $|\gamma|\ge s-j$, and
\[
  x^\alpha\in\mm_x^{s-j}Q_x^{\,j}\iff \fl{\alpha/2}\ge j\ \text{and}\ |\alpha|\ge s+j.
\]
Indeed, $\mm_x^{s-j}Q_x^{\,j}$ is generated by the monomials
$x^{2\beta+\delta}$ with $|\beta|=j$ and $|\delta|=s-j$. If such a monomial
divides $x^\alpha$, then $\beta_a\le\fl{\alpha_a/2}$ for all $a$, and
$|\alpha|\ge2j+(s-j)$. Conversely, if $\fl{\alpha/2}\ge j$ and $|\alpha|\ge s+j$,
choose $\beta$ with $|\beta|=j$ and $\beta_a\le\fl{\alpha_a/2}$; then
$\alpha-2\beta\in\NN^p$ has degree $|\alpha|-2j\ge s-j$ and contains some
$\delta$ with $|\delta|=s-j$.

Hence $x^\alpha y^\gamma\in I^s$ if and only if one of the following holds for
some $0\le j\le s$:
\begin{enumerate}[label=(\roman*)]
\item $\fl{\alpha/2}\ge j$, $|\alpha|\ge s+j$ and $|\gamma|\ge s-j$;
\item $\fl{\gamma/2}\ge j$, $|\gamma|\ge s+j$ and $|\alpha|\ge s-j$.
\end{enumerate}
We claim that (i) holds for some $j$ if and only if
\begin{equation}\label{eq:condX}
  |\alpha|\ge s,\qquad |\alpha|+|\gamma|\ge2s,\qquad \fl{\alpha/2}+|\gamma|\ge s.
\end{equation}
If (i) holds, then $|\alpha|\ge s+j\ge s$, $|\alpha|+|\gamma|\ge2s$ and
$\fl{\alpha/2}+|\gamma|\ge j+(s-j)=s$. Conversely, assume~\eqref{eq:condX} and
take $j=\max\{0,s-|\gamma|\}\le s$. If $j=0$, condition (i) reduces to
$|\alpha|\ge s$. If $j=s-|\gamma|>0$, then $\fl{\alpha/2}\ge s-|\gamma|=j$ and
$|\alpha|\ge2s-|\gamma|=s+j$, and $|\gamma|=s-j$. This proves the claim, and
symmetrically (ii) holds for some $j$ if and only if $|\gamma|\ge s$,
$|\alpha|+|\gamma|\ge2s$ and $\fl{\gamma/2}+|\alpha|\ge s$.

If~\eqref{eq:condX} holds, then $\fl{\gamma/2}+|\alpha|\ge|\alpha|\ge s$, so the
three conditions of the lemma hold, and symmetrically for (ii). Conversely, if
the three conditions hold, then $|\alpha|\ge s$ or $|\gamma|\ge s$, and
accordingly~\eqref{eq:condX} or its symmetric counterpart holds.
\end{proof}

\begin{remark}\label{rem:order}
Equivalently, the largest $s$ with $x^\alpha y^\gamma\in I_{p,q}^s$ is
\[
  \min\Bigl\{\Bigl\lfloor\frac{|\alpha|+|\gamma|}{2}\Bigr\rfloor,\
  \fl{\alpha/2}+|\gamma|,\ \fl{\gamma/2}+|\alpha|\Bigr\}.
\]
The first quantity is the bound coming from degrees. The second reflects the
absence of the products $x_ax_{a'}$, $a\ne a'$, from $I_{p,q}$: in a product of
generators dividing $x^\alpha y^\gamma$, at most $\fl{\alpha/2}$ factors are
squares of $x$-variables, and every other factor involves a $y$-variable. The
third is symmetric.
\end{remark}

\section{Lengths of powers}\label{sec:lengths}

For integers $a,b\ge1$ we let $D_{a,b}(s)$ be as in~\eqref{eq:D}. Note that
$D_{1,b}=0$, because the inner sum requires $t+2r\le a$.

\begin{proposition}\label{prop:length}
For every $s\ge1$,
\[
  \ell(S/I_{p,q}^s)=\binom{2s+n-1}{n}+D_{p,q}(s)+D_{q,p}(s).
\]
\end{proposition}

\begin{proof}
Write $I=I_{p,q}$. We count the standard monomials of $I^s$ by
means of Lemma~\ref{lem:membership}. Every monomial of degree at most $2s-1$ is
standard, and there are $\binom{2s-1+n}{n}$ of them. A monomial
$x^\alpha y^\gamma$ of degree at least $2s$ is standard if and only if
$\fl{\alpha/2}+|\gamma|<s$ or $\fl{\gamma/2}+|\alpha|<s$. The two inequalities
cannot hold simultaneously, because the sum of their left-hand sides is at least
$|\alpha|+|\gamma|\ge2s$.

We count the monomials of degree at least $2s$ with
$\fl{\alpha/2}+|\gamma|<s$. Write $\alpha=2\beta+\varepsilon$ with
$\beta\in\NN^p$ and $\varepsilon\in\{0,1\}^p$; this is a bijection between
$\NN^p$ and such pairs, and $|\varepsilon|=\omega(\alpha)$,
$|\beta|=\fl{\alpha/2}$. Put $t=|\gamma|$, $\kappa=|\varepsilon|$, $l=|\beta|$ and
$r=s-t-l$. The condition $\fl{\alpha/2}+|\gamma|<s$ says $r\ge1$, and the degree
condition $2l+\kappa+t\ge2s=2t+2l+2r$ says $\kappa\ge t+2r$. Conversely, given $t\ge0$,
$r\ge1$ with $t+r\le s$ and $t+2r\le \kappa\le p$, there are $\binom{t+q-1}{q-1}$
choices of $\gamma$, $\binom{p}{\kappa}$ choices of $\varepsilon$ and
$\binom{s-t-r+p-1}{p-1}$ choices of $\beta$ with $|\beta|=s-t-r$. Hence these
monomials are counted by $D_{p,q}(s)$. By symmetry the monomials of degree at
least $2s$ with $\fl{\gamma/2}+|\alpha|<s$ are counted by $D_{q,p}(s)$.
\end{proof}

For $s=1$ this gives $\ell(S/I_{p,q})=2^p+2^q-1$, the standard monomials being
the squarefree monomials in the $x$-variables and those in the $y$-variables.

For integers $a,b\ge1$ define the polynomial
\begin{equation}\label{eq:P}
  P_{a,b}(z)=\sum_{\substack{t\ge0,\ r\ge1\\ t+2r\le a}}
  \binom{t+b-1}{b-1}\Biggl(\,\sum_{\kappa=t+2r}^{a}\binom{a}{\kappa}\Biggr)z^{t+r-1}.
\end{equation}
Thus $P_{1,b}=0$, and for $a\ge2$ the degree of $P_{a,b}$ is at most $a-2$,
because $t+r-1\le t+2r-2\le a-2$. For $a=0$ the sum in~\eqref{eq:P} is
empty, and we put $P_{0,b}=0$.

\begin{lemma}\label{lem:Dseries}
For $a,b\ge1$,
\[
  \sum_{s\ge1}D_{a,b}(s)z^s=\frac{z\,P_{a,b}(z)}{(1-z)^a}.
\]
\end{lemma}

\begin{proof}
Substitute $l=s-t-r\ge0$ in~\eqref{eq:D} and use
$\sum_{l\ge0}\binom{l+a-1}{a-1}z^l=(1-z)^{-a}$.
\end{proof}

The main term of Proposition~\ref{prop:length} is the length of $S/\mm^{2s}$,
where $\mm=\mm_x+\mm_y$. Its Hilbert coefficients are as follows.

\begin{proposition}\label{prop:square}
Let $n\ge1$ and $h_n(z)=\sum_{j\ge0}\binom{n+1}{2j+1}z^j$. Then
\begin{enumerate}[label=\textup{(\roman*)}]
\item $\displaystyle\sum_{s\ge1}\binom{2s+n-1}{n}z^s=\frac{z\,h_n(z)}{(1-z)^{n+1}}$;
\item $\displaystyle h_n(1+u)=\sum_{i=0}^{\fl{n/2}}2^{n-2i}\binom{n-i}{i}u^i$.
\end{enumerate}
Consequently, if $\mm$ is the homogeneous maximal ideal of a polynomial ring in
$n$ variables, then $e_i(\mm^2)=2^{n-2i}\binom{n-i}{i}$ for $0\le i\le\fl{n/2}$,
$e_i(\mm^2)=0$ for $\fl{n/2}<i\le n$, and
\[
  \binom{2s+n-1}{n}=\sum_{i=0}^{\fl{n/2}}(-1)^i2^{n-2i}\binom{n-i}{i}\binom{s+n-i-1}{n-i}
  \qquad\text{for every }s\ge1.
\]
\end{proposition}

\begin{proof}
(i) From $\sum_{m\ge0}\binom{m+n}{n}w^m=(1-w)^{-n-1}$ we obtain
\begin{align*}
  \sum_{s\ge1}\binom{2s-1+n}{n}w^{2s-1}
  &=\frac{(1-w)^{-n-1}-(1+w)^{-n-1}}{2}\\
  &=\frac{(1+w)^{n+1}-(1-w)^{n+1}}{2(1-w^2)^{n+1}}\\
  &=\frac{1}{(1-w^2)^{n+1}}\sum_{j\ge0}\binom{n+1}{2j+1}w^{2j+1}.
\end{align*}
Multiplying by $w$ and writing $z=w^2$ gives (i).

(ii) For $c\ge1$ we have $\sum_{n\ge0}\binom{n+1}{c}v^n=v^{c-1}/(1-v)^{c+1}$.
Hence, in $\QQ[u][[v]]$,
\begin{align*}
  \sum_{n\ge0}h_n(1+u)v^n&=\sum_{j\ge0}\frac{(1+u)^jv^{2j}}{(1-v)^{2j+2}}
  =\frac{1}{(1-v)^2-(1+u)v^2}\\
  &=\frac{1}{1-v(2+uv)}=\sum_{m\ge0}v^m(2+uv)^m,
\end{align*}
and the coefficient of $u^iv^n$ on the right is the coefficient of $u^i$ in
$(2+u)^{n-i}$, namely $\binom{n-i}{i}2^{n-2i}$ if $2i\le n$ and $0$ otherwise.

Since $\ell(S/(\mm^2)^s)=\binom{2s+n-1}{n}$ and $\deg h_n=\fl{n/2}\le n$, the
last assertions follow from Lemma~\ref{lem:numerator}.
\end{proof}

Proposition~\ref{prop:square} proves the formula
$e_j=\binom{n-j}{j}2^{n-2j}$ conjectured in~\cite[Conjecture~3.1]{Fr}, including
the assertion that $e_i(\mm^2)\ne0$ exactly when $i\le\fl{n/2}$. The binomial
identity stated there differs from the last display of
Proposition~\ref{prop:square} in its last binomial coefficient, which is
printed as $\binom{s+n-j}{n-j}$; by~\eqref{eq:convention} it should be
$\binom{s+n-j-1}{n-j}$. For $n=2$ and $s=1$ the printed right-hand side equals
$10$, whereas the left-hand side is $3$. The numerical
example~\cite[Example~3.1]{Fr} is in the corrected form.

\begin{corollary}\label{cor:numerator}
Let
\[
  h_{p,q}(z)=h_n(z)+(1-z)^{q+1}P_{p,q}(z)+(1-z)^{p+1}P_{q,p}(z).
\]
Then $\sum_{s\ge1}\ell(S/I_{p,q}^s)z^s=z\,h_{p,q}(z)/(1-z)^{n+1}$, and
$\deg h_{p,q}\le n-1$. Consequently~\eqref{eq:convention} holds for every
$s\ge1$, with
\[
  e_i(I_{p,q})=\cf{u^i}h_{p,q}(1+u)\qquad(0\le i\le n).
\]
Moreover $e_n(I_{p,q})=0$, $e_0(I_{p,q})=2^n$, and
$e_i(I_{p,q})=2^{n-2i}\binom{n-i}{i}>0$ for $0\le i\le p$.
\end{corollary}

\begin{proof}
By Proposition~\ref{prop:length}, Lemma~\ref{lem:Dseries} and
Proposition~\ref{prop:square}(i), the generating function of the lengths is
\[
  \frac{z\,h_n(z)}{(1-z)^{n+1}}+\frac{z\,P_{p,q}(z)}{(1-z)^{p}}
  +\frac{z\,P_{q,p}(z)}{(1-z)^{q}}
  =\frac{z\,h_{p,q}(z)}{(1-z)^{n+1}}.
\]
Since $n\ge2$, $\deg h_n=\fl{n/2}\le n-1$, and
$\deg\,(1-z)^{q+1}P_{p,q}\le q+1+p-2=n-1$; the same bound holds for the third
summand. Lemma~\ref{lem:numerator} gives~\eqref{eq:convention} for all $s\ge1$,
the formula for $e_i$, and $e_n=0$. At $z=1+u$ the second and third summands of
$h_{p,q}$ become $(-u)^{q+1}P_{p,q}(1+u)$ and $(-u)^{p+1}P_{q,p}(1+u)$, which are
divisible by $u^{p+1}$. Hence $e_i(I_{p,q})=e_i(\mm^2)$ for $i\le p$, and
Proposition~\ref{prop:square} applies. Finally $\binom{n-i}{i}>0$ for
$2i\le 2p\le n$.
\end{proof}

This proves Theorem~\ref{thm:main}(a) and the statements of
Theorem~\ref{thm:main}(b) that follow the words ``in particular''.

\section{The coefficient formula}\label{sec:coefficients}

By Corollary~\ref{cor:numerator} it remains to compute
$\cf{u^i}(-u)^{a+1}P_{b,a}(1+u)$ for $(a,b)=(q,p)$ and $(a,b)=(p,q)$. We first
sum the polynomials $P_{b,a}$ over $b$.

\begin{lemma}\label{lem:Pseries}
Let $a\ge1$. In $\QQ[z][[y]]$,
\[
  \sum_{b\ge0}P_{b,a}(z)\,y^b
  =\frac{y^2(1-y)^a}{(1-2y)\,\bigl(1-(1+z)y\bigr)^a\,\bigl((1-y)^2-zy^2\bigr)}.
\]
\end{lemma}

\begin{proof}
For $\kappa\ge0$ let
$Q_\kappa(z)=\sum_{t\ge0,\,r\ge1,\,t+2r\le \kappa}\binom{t+a-1}{a-1}z^{t+r-1}$;
in particular $Q_0=0$. Exchanging the summations over $\kappa$ and over $(t,r)$
in~\eqref{eq:P} gives $P_{b,a}=\sum_{\kappa=0}^{b}\binom{b}{\kappa}Q_\kappa$ for
all $b\ge0$; for $b=0$ both sides vanish, because $P_{0,a}=0$ and $Q_0=0$.
Summing over $t\ge0$, $r\ge1$ and $\kappa\ge t+2r$ independently,
\[
  \sum_{\kappa\ge0}Q_\kappa(z)w^\kappa=\frac{1}{1-w}\cdot\frac{1}{(1-zw)^a}\cdot\frac{w^2}{1-zw^2}.
\]
Since $\sum_{b\ge \kappa}\binom{b}{\kappa}y^b=y^\kappa/(1-y)^{\kappa+1}$, we get
\[
  \sum_{b\ge0}P_{b,a}(z)\,y^b=\sum_{\kappa\ge0}Q_\kappa(z)\frac{y^\kappa}{(1-y)^{\kappa+1}}
  =\frac{1}{1-y}\sum_{\kappa\ge0}Q_\kappa(z)w^\kappa\Big|_{w=y/(1-y)}.
\]
With $w=y/(1-y)$ we have $1-w=(1-2y)/(1-y)$, $1-zw=(1-(1+z)y)/(1-y)$ and
$1-zw^2=((1-y)^2-zy^2)/(1-y)^2$, and the stated formula follows.
\end{proof}

\begin{proposition}\label{prop:correction}
Let $a,b\ge1$ and $i\ge1$. Then
\[
  \cf{u^i}\Bigl((-u)^{a+1}P_{b,a}(1+u)\Bigr)
  =-\cf{z^{a+b-i-1}}\frac{\Delta_{a,i}(z)}{(1-2z)^{i+1}}.
\]
\end{proposition}

\begin{proof}
If $i\le a$, both sides vanish: the left-hand side because of the factor
$u^{a+1}$, the right-hand side because $\Delta_{a,i}=0$. Let $i\ge a+1$. After
multiplication by $(-1)^{a+1}$, the claim reads
\begin{equation}\label{eq:claim-b}
  \cf{u^{i-a-1}}P_{b,a}(1+u)
  =(-1)^a\,\cf{z^{a+b-i-1}}\frac{\Delta_{a,i}(z)}{(1-2z)^{i+1}}.
\end{equation}
We prove~\eqref{eq:claim-b} for all $b\ge0$, where $P_{0,a}=0$. The right-hand side vanishes for
$b\le i-a$, and its generating function over $b\ge0$ is
$(-1)^ay^{i+1-a}\Delta_{a,i}(y)/(1-2y)^{i+1}$. By Lemma~\ref{lem:Pseries} with
$z=1+u$, the generating function of the left-hand side over $b\ge0$ is
$\cf{u^{i-a-1}}\Phi$, where
\[
  \Phi=\Phi(u,y)=\frac{y^2(1-y)^a}{(1-2y)\,\bigl(1-(2+u)y\bigr)^a\,(1-2y-uy^2)}
  \in\QQ[u][[y]].
\]
Comparing coefficients of the powers of $u$, we see that~\eqref{eq:claim-b}
holds for all $b\ge0$ and all $i\ge a+1$ if and only if
\begin{equation}\label{eq:claim-sum}
  \sum_{i\ge a+1}u^{i-a-1}\,\frac{y^{i+1-a}\,\Delta_{a,i}(y)}{(1-2y)^{i+1}}
  =(-1)^a\,\Phi
\end{equation}
in $\QQ[u][[y]]$; the series on the left converges $y$-adically, since its
$i$-th term is divisible by $y^{i+1-a}$. Put $w=uy/(1-2y)\in y\QQ[u][[y]]$. The
$i$-th term on the left of~\eqref{eq:claim-sum} equals
$y^2(1-2y)^{-a-2}\,w^{i-a-1}\Delta_{a,i}(y)$, and with $m=i-1$,
\begin{align*}
  \sum_{i\ge a+1}w^{i-a-1}\Delta_{a,i}(y)
  &=\sum_{j\ge a}(-1)^j(1-y)^j\sum_{m\ge j}\binom{m}{j}w^{m-a}\\
  &=\sum_{j\ge a}(-1)^j(1-y)^j\frac{w^{j-a}}{(1-w)^{j+1}}\\
  &=\frac{(-1)^a(1-y)^a}{(1-w)^{a+1}}\cdot\frac{1}{1+(1-y)w/(1-w)}\\
  &=\frac{(-1)^a(1-y)^a}{(1-w)^a(1-yw)}.
\end{align*}
Finally $(1-2y)^a(1-w)^a=(1-(2+u)y)^a$ and $(1-2y)(1-yw)=1-2y-uy^2$, so that
\[
  \frac{y^2}{(1-2y)^{a+2}}\cdot\frac{(-1)^a(1-y)^a}{(1-w)^a(1-yw)}=(-1)^a\Phi.
\]
This proves~\eqref{eq:claim-sum} and the proposition.
\end{proof}

\begin{proof}[Proof of Theorem~\ref{thm:main}\textup{(b)}]
By Corollary~\ref{cor:numerator} it remains to prove the formula for $e_i$,
$1\le i\le n$. By Corollary~\ref{cor:numerator} and
Proposition~\ref{prop:square}(ii),
\[
  e_i(I_{p,q})=2^{n-2i}\binom{n-i}{i}
  +\cf{u^i}(-u)^{q+1}P_{p,q}(1+u)+\cf{u^i}(-u)^{p+1}P_{q,p}(1+u).
\]
The first term equals
$\cf{z^{n-2i}}(1-2z)^{-i-1}=\cf{z^{n-i-1}}z^{i-1}/(1-2z)^{i+1}$, and both sides
vanish if $2i>n$. By Proposition~\ref{prop:correction} with $(a,b)=(q,p)$ and
$(a,b)=(p,q)$, the second and third terms are
$-\cf{z^{n-i-1}}\Delta_{q,i}(z)/(1-2z)^{i+1}$ and
$-\cf{z^{n-i-1}}\Delta_{p,i}(z)/(1-2z)^{i+1}$.
\end{proof}

\section{Signs}\label{sec:signs}

For $i\ge2$ and $1\le m\le i-1$ put
\[
  G_{m,i}(z)=z^{i-1}-\Delta_{m,i}(z)=\sum_{j=0}^{m-1}(-1)^j\binom{i-1}{j}(1-z)^j,
\]
where the second expression follows from~\eqref{eq:Delta} and the binomial
expansion $z^{i-1}=\bigl(1-(1-z)\bigr)^{i-1}=\sum_{j=0}^{i-1}(-1)^j\binom{i-1}{j}(1-z)^j$.
For $r\ge0$ put
\[
  \tau_{m,i}(r)=\cf{z^r}\frac{G_{m,i}(z)}{(1-2z)^{i+1}}.
\]

\begin{lemma}\label{lem:reduction}
Let $p+1\le i\le n-1$ and $r=n-i-1\ge0$. Then
\[
  e_i(I_{p,q})=
  \begin{cases}
    \tau_{p,i}(r) & \text{if } i\le q,\\
    \tau_{p,i}(r)+\tau_{q,i}(r) & \text{if } i>q.
  \end{cases}
\]
\end{lemma}

\begin{proof}
Since $p<i$, the numerator in Theorem~\ref{thm:main}(b) is
$G_{p,i}-\Delta_{q,i}$. If $i\le q$ then $\Delta_{q,i}=0$. If $i>q$, the
numerator is $G_{p,i}+G_{q,i}-z^{i-1}$; since $n=p+q\le2i-2$, we have
$r\le i-3$, so $z^{i-1}$ does not contribute to the coefficient of $z^r$.
\end{proof}

\begin{lemma}\label{lem:partial}
For integers $M\ge0$ and $N\ge M+1$ put
\[
  \beta_M(N)=(-1)^M\sum_{l=0}^{M}\Bigl(-\frac12\Bigr)^l\binom{N}{l}.
\]
Then $\beta_M(N)\ge0$, and $\beta_0(N)=1$.
\end{lemma}

\begin{proof}
The value $\beta_0(N)=1$ is clear. We prove $\beta_M(N)\ge0$ by induction on
$N\ge1$, simultaneously for all $M$ with $0\le M\le N-1$. The case $M=0$ is
clear. For $N=M+1$,
\[
  \sum_{l=0}^{M}\Bigl(-\frac12\Bigr)^l\binom{M+1}{l}
  =\Bigl(1-\frac12\Bigr)^{M+1}-\Bigl(-\frac12\Bigr)^{M+1},
\]
so $\beta_M(M+1)=2^{-M-1}\bigl(1+(-1)^M\bigr)\ge0$. This settles $N=1$. For
$1\le M\le N-2$, Pascal's rule gives
\[
  \beta_M(N)=\beta_M(N-1)+\tfrac12\,\beta_{M-1}(N-1),
\]
and both terms on the right are nonnegative by the induction hypothesis for
$N-1$, since $M\le (N-1)-1$ and $M-1\le(N-1)-1$.
\end{proof}

\begin{lemma}\label{lem:positive}
Let $i\ge2$ and $1\le m\le i-1$. Then
\[
  (-1)^{m-1}G_{m,i}(z)=\sum_{c=0}^{m-1}\gamma_c\,(1-2z)^c
\]
with rational numbers $\gamma_c\ge0$ and $\gamma_{m-1}>0$. Consequently
$(-1)^{m-1}\tau_{m,i}(r)>0$ for every $r\ge0$.
\end{lemma}

\begin{proof}
Write $v=1-2z$, so that $1-z=(1+v)/2$. The coefficient of $v^c$ in $G_{m,i}$ is
\[
  \sum_{j=c}^{m-1}(-1)^j2^{-j}\binom{i-1}{j}\binom{j}{c}
  =(-1)^c2^{-c}\binom{i-1}{c}\sum_{l=0}^{m-1-c}\Bigl(-\frac12\Bigr)^l\binom{i-1-c}{l},
\]
using $\binom{i-1}{j}\binom{j}{c}=\binom{i-1}{c}\binom{i-1-c}{j-c}$ and $l=j-c$.
Put $M=m-1-c$ and $N=i-1-c\ge M+1$. Since $(-1)^c(-1)^M=(-1)^{m-1}$, the
coefficient of $v^c$ in $(-1)^{m-1}G_{m,i}$ is
$\gamma_c=2^{-c}\binom{i-1}{c}\beta_M(N)$, which is nonnegative by
Lemma~\ref{lem:partial}, and $\gamma_{m-1}=2^{1-m}\binom{i-1}{m-1}>0$. Hence
\[
  (-1)^{m-1}\tau_{m,i}(r)=\sum_{c=0}^{m-1}\gamma_c\,\cf{z^r}(1-2z)^{c-i-1}
  =\sum_{c=0}^{m-1}\gamma_c\,2^r\binom{r+i-c}{r}>0,
\]
since $c-i-1<0$ for all $c\le m-1$.
\end{proof}

Lemmas~\ref{lem:reduction} and~\ref{lem:positive} already give the sign of
$e_i(I_{p,q})$ when $i\le q$, and when $i>q$ and $p\equiv q\pmod 2$. For $i>q$
and $p\not\equiv q\pmod 2$ the two contributions have opposite signs, and we
show that the contribution of the smaller part dominates. We use the following
integral representation; here $\int_0^1$ denotes the $\QQ$-linear map on
polynomials in $\theta$ that sends $\theta^c$ to $1/(c+1)$, applied coefficientwise.

\begin{lemma}\label{lem:integral}
Let $i\ge2$ and $1\le m\le i-1$. In $\QQ[z]$,
\[
  G_{m,i}(z)-z^{i-1}=(-1)^{m-1}(i-1)\binom{i-2}{m-1}(1-z)^m
  \int_0^1 \theta^{m-1}\bigl(z+(1-z)\theta\bigr)^{i-m-1}\,d\theta.
\]
\end{lemma}

\begin{proof}
For integers $N\ge m\ge1$ we have in $\QQ[X,Y]$
\begin{equation}\label{eq:beta}
  \sum_{j=m}^{N}\binom{N}{j}X^jY^{N-j}
  =N\binom{N-1}{m-1}X^m\int_0^1 \theta^{m-1}\bigl(Y+X(1-\theta)\bigr)^{N-m}\,d\theta.
\end{equation}
Indeed, by $\int_0^1\theta^{m-1}(1-\theta)^l\,d\theta=(m-1)!\,l!/(m+l)!$ the coefficient of
$X^{m+l}Y^{N-m-l}$ on the right is
$N\binom{N-1}{m-1}\binom{N-m}{l}\frac{(m-1)!\,l!}{(m+l)!}=\binom{N}{m+l}$. Now
$G_{m,i}(z)-z^{i-1}=-\Delta_{m,i}(z)=-\sum_{j=m}^{i-1}\binom{i-1}{j}(z-1)^j$.
Apply~\eqref{eq:beta} with $N=i-1$, $X=z-1$ and $Y=1$, and note
$Y+X(1-\theta)=z+(1-z)\theta$ and $X^m=(-1)^m(1-z)^m$.
\end{proof}

\begin{lemma}\label{lem:substitution}
Let $f\in\QQ[[z]]$ and $r\ge0$. Then
$\cf{z^r}f(z)=\cf{w^r}(1+2w)^{r-1}f\bigl(w/(1+2w)\bigr)$.
\end{lemma}

\begin{proof}
Both sides depend only on the coefficients of $z^0,\dots,z^r$ in $f$ and are
linear in $f$, so it suffices to take $f=z^c$ with $0\le c\le r$. The right-hand
side is then $\cf{w^{r-c}}(1+2w)^{r-1-c}$, which is $1$ if $c=r$ and $0$ if
$c<r$, since then $(1+2w)^{r-1-c}$ is a polynomial of degree $r-1-c$.
\end{proof}

\begin{proposition}\label{prop:monotone}
Let $i\ge2$ and $0\le r\le i-2$. For $1\le m\le i-1$ put
\[
  \upsilon_m=\frac{(-1)^{m-1}\tau_{m,i}(r)}{\binom{i-2}{m-1}}.
\]
Then $\upsilon_1\ge\upsilon_2\ge\cdots\ge\upsilon_{i-1}>0$.
\end{proposition}

\begin{proof}
Positivity is Lemma~\ref{lem:positive}. Since $r<i-1$, the coefficient of $z^r$
in $z^{i-1}/(1-2z)^{i+1}$ vanishes, so by Lemma~\ref{lem:integral}
\[
  \upsilon_m=(i-1)\,\cf{z^r}\frac{(1-z)^m}{(1-2z)^{i+1}}
  \int_0^1 \theta^{m-1}\bigl(z+(1-z)\theta\bigr)^{i-m-1}\,d\theta.
\]
We apply Lemma~\ref{lem:substitution}. For $z=w/(1+2w)$ we have
$1-2z=(1+2w)^{-1}$, $1-z=(1+w)/(1+2w)$ and $z+(1-z)\theta=(\theta+(1+\theta)w)/(1+2w)$. Since
evaluation at $z=w/(1+2w)$ commutes with $\int_0^1$, and since
$(r-1)+(i+1)-m-(i-m-1)=r+1$, we obtain
\[
  \upsilon_m=(i-1)\,\cf{w^r}(1+2w)^{r+1}\int_0^1F_m(\theta,w)\,d\theta,
\]
where $F_m=(1+w)^m\theta^{m-1}\bigl(\theta+(1+\theta)w\bigr)^{i-m-1}$.
For $1\le m\le i-2$,
\begin{align*}
  F_m-F_{m+1}&=(1+w)^m\theta^{m-1}\bigl(\theta+(1+\theta)w\bigr)^{i-m-2}
  \bigl(\theta+(1+\theta)w-(1+w)\theta\bigr)\\
  &=w\,(1+w)^m\theta^{m-1}\bigl(\theta+(1+\theta)w\bigr)^{i-m-2},
\end{align*}
a polynomial in $\theta$ and $w$ with nonnegative coefficients. Multiplication by
$(1+2w)^{r+1}$, the map $\int_0^1$ and the extraction of a coefficient preserve
nonnegativity, so $\upsilon_m-\upsilon_{m+1}\ge0$.
\end{proof}

\begin{proof}[Proof of Theorem~\ref{thm:main}\textup{(c)} and~\textup{(d)}]
Let $p+1\le i\le n-1$ and $r=n-i-1\ge0$. By Lemma~\ref{lem:positive},
$\tau_{p,i}(r)$ has sign $(-1)^{p-1}=(-1)^{p+1}$, and if $q<i$ then
$\tau_{q,i}(r)$ has sign $(-1)^{q+1}$. By Lemma~\ref{lem:reduction}, this proves
(c) if $i\le q$, and also if $i>q$ and $p\equiv q\pmod2$.

Let $i>q$ and $p\not\equiv q\pmod 2$, so $p<q<i$. As in
Lemma~\ref{lem:reduction}, $r\le i-3$, so Proposition~\ref{prop:monotone}
applies with $m\in\{p,q\}$ and gives $\upsilon_p\ge\upsilon_q>0$. We claim
$\binom{i-2}{p-1}>\binom{i-2}{q-1}$. Put $N=i-2$ and $c=N-(q-1)\ge0$. Then
$\binom{N}{q-1}=\binom{N}{c}$, and $c<p-1<N-c$, because $(p-1)+(q-1)=n-2>N$ and
$p<q$. Since $\min\{p-1,N-p+1\}>c$ and the binomial coefficients
$\binom{N}{j}$ increase strictly for $0\le j\le\fl{N/2}$, the claim follows.
Hence
\[
  (-1)^{p+1}\tau_{p,i}(r)=\binom{i-2}{p-1}\upsilon_p
  >\binom{i-2}{q-1}\upsilon_q=(-1)^{q+1}\tau_{q,i}(r)>0,
\]
and $e_i(I_{p,q})=\tau_{p,i}(r)+\tau_{q,i}(r)$ has the sign $(-1)^{p+1}$. This
proves~(c).

For (d), if $p$ is odd then $e_0,\dots,e_{n-1}$ are positive by (b) and (c), and
$e_n=0$. If $p$ is even, then $p\ge2$, so $q\ge2$ and $n-1\ge p+1$; by (c) the
$q-1$ coefficients $e_{p+1},\dots,e_{n-1}$ are negative.
\end{proof}

\section{Small part sizes and comparison with the literature}\label{sec:examples}

\begin{corollary}[Stars]\label{cor:star}
Let $q\ge1$ and $n=q+1$. Then $e_0(I_{1,q})=2^{q+1}$,
$e_i(I_{1,q})=2^{q-i}\binom{q}{i}$ for $1\le i\le q$, and $e_{q+1}(I_{1,q})=0$.
Equivalently, for every $s\ge1$,
\[
  \ell(S/I_{1,q}^s)=2^{q+1}\binom{s+q}{q+1}
  +\sum_{i=1}^{q}(-1)^i2^{q-i}\binom{q}{i}\binom{s+q-i}{q+1-i}.
\]
\end{corollary}

\begin{proof}
For $i=1$ Theorem~\ref{thm:main}(b) gives $2^{n-2}\binom{n-1}{1}=2^{q-1}q$. For
$2\le i\le q$ we have $i\le q$ and $G_{1,i}=1$, so Lemma~\ref{lem:reduction}
gives $e_i=\cf{z^{q-i}}(1-2z)^{-i-1}=2^{q-i}\binom{q}{i}$.
\end{proof}

For $q=2$ this is the case $M=\{x_1x_2,x_1x_3\}$ of \cite[Theorem~3.1]{Fr},
with $(e_0,e_1,e_2)=(8,4,1)$.

\begin{proof}[Proof of Corollary~\ref{cor:k2q}]
The values of $e_0,e_1,e_2$ and $e_n$ are Theorem~\ref{thm:main}(b). Let
$3\le i\le n-1$ and $r=n-i-1$. Here $G_{2,i}(z)=1-(i-1)(1-z)=(2-i)+(i-1)z$, so
\[
  \tau_{2,i}(r)=(2-i)\,2^r\binom{r+i}{i}+(i-1)\,2^{r-1}\binom{r+i-1}{i}
\]
for $r\ge1$, and $\tau_{2,i}(0)=2-i$. If $i\le n-2=q$, then
$e_i=\tau_{2,i}(r)$ by Lemma~\ref{lem:reduction}. Using
$\binom{r+i}{i}=\frac{r+i}{i}\binom{r+i-1}{i-1}$ and
$\binom{r+i-1}{i}=\frac{r}{i}\binom{r+i-1}{i-1}$ we get
\begin{align*}
  \tau_{2,i}(r)&=\frac{2^{r-1}}{i}\binom{r+i-1}{i-1}\bigl(2(2-i)(r+i)+(i-1)r\bigr)\\
  &=-\frac{2^{r-1}}{i}\binom{r+i-1}{i-1}\bigl(2i(i-2)+(i-3)r\bigr),
\end{align*}
which is negative for $i\ge3$ and $r\ge1$. If $i=n-1=q+1$, then $r=0$ and, by
Lemma~\ref{lem:reduction}, $e_i=\tau_{2,i}(0)+\tau_{q,i}(0)$, where
$\tau_{q,i}(0)=G_{i-1,i}(0)=\sum_{j=0}^{i-2}(-1)^j\binom{i-1}{j}=(-1)^i$. Thus
$e_{n-1}=2-(n-1)+(-1)^{n-1}=3-n+(-1)^{n-1}<0$.
\end{proof}

\begin{proof}[Proof of Corollary~\ref{cor:c4}]
For $p=q=2$ the only term in~\eqref{eq:D} is $t=0$, $r=1$, $\kappa=2$, so
$D_{2,2}(s)=s$, and Theorem~\ref{thm:main}(a) gives
$\ell(S/I_{2,2}^s)=\binom{2s+3}{4}+2s$. The coefficients are those of
Corollary~\ref{cor:k2q} with $n=4$, and the displayed expansion
is~\eqref{eq:convention}. Since $S/I_{2,2}$ is an Artinian quadratic algebra
and $e_3(I_{2,2})=-2$, the ideal $I_{2,2}$ contradicts
\cite[Conjecture~1.1]{Fr}.
\end{proof}

\begin{remark}\label{rem:froberg}
The four-variable predictions of \cite[Conjecture~3.2]{Fr} that belong to the
family $I_{p,q}$ are the cases $p=q=2$ and $(p,q)=(1,3)$. In both cases the
predicted numbers of standard monomials of degree at least $2s$, namely $2s$
for $M=\{x_1x_2,x_2x_3,x_3x_4,x_4x_1\}$ and $s(2s+1)+\binom{s+1}{2}$ for
$M=\{x_1x_2,x_1x_3,x_1x_4\}$, agree with $D_{p,q}(s)+D_{q,p}(s)$ in
Theorem~\ref{thm:main}(a). For the four-cycle the listed coefficient $e_3=2$
should be $e_3=-2$, as explained after Corollary~\ref{cor:c4}. For the star
$K_{1,3}$, the listed values $(e_0,e_1,e_2)=(16,12,6)$ agree with
Corollary~\ref{cor:star}, which gives in addition $e_3=1$ and $e_4=0$.

The case $M=\{x_1x_2,x_1x_3,x_1x_4,x_2x_3,x_2x_4\}$, the complete graph on four
vertices with the edge $x_3x_4$ removed, is not in the family. Its predicted
length is $\binom{2s+3}{4}+s$. Under~\eqref{eq:convention} this gives
$(e_0,\dots,e_4)=(16,12,1,-1,0)$, whereas the listed coefficients are
$(16,12,1,1)$.
\end{remark}

\begin{remark}\label{rem:bgv}
Banerjee, Goel and Verma~\cite[Theorem~3.3, Example~3.11]{BGV} computed the
lengths $\ell\bigl(R/(\mathfrak n^{[Q]})^t\bigr)$ for the Stanley--Reisner
ring $R=S/I(K_{\alpha,\beta})$ defined by the edge ideal of a complete bipartite
graph, with $1\le\alpha\le\beta$ as in~\cite{BGV},
where $\mathfrak n$ is the homogeneous maximal ideal of $R$ and $\mathfrak
n^{[Q]}$ is generated by the $Q$-th powers of the variables. There the edge
ideal is fixed and only $\mathfrak n^{[Q]}$ is raised to the $t$-th power, and
the ring $R$ has dimension $\beta$, so that the corresponding Hilbert--Samuel
polynomial in $t$ has degree $\beta$. In Theorem~\ref{thm:main} the whole ideal
$I(K_{p,q})+(x_a^2,y_b^2)$ is raised to the $s$-th power in the polynomial ring
$S$ of dimension $n$, and the polynomial in $s$ has degree $n$. For $Q=2$ the
two lengths agree when $s=t=1$, both being $2^p+2^q-1$, but not in general: for
$p=q=2$ and $s=t=2$ the formula
of~\cite[Example~3.11]{BGV} gives $23$, whereas $\ell(S/I_{2,2}^2)=39$.
\end{remark}

\begin{remark}\label{rem:computations}
The results were checked by exact computer calculations, which are not used in
the proofs. Lemma~\ref{lem:membership} was compared with a direct recursive
computation of the largest power of $I_{p,q}$ containing a given monomial, for
all monomials of degree at most $9$ with $p\le4$ and $p\le q\le5$, and
Theorem~\ref{thm:main}(a) with a direct count of the standard monomials of
$I_{p,q}^s$ for $p\le2$, $q\le4$ and $s\le3$; both computations start from the
generators of $I_{p,q}$ and are independent of our arguments. For $p\le16$ and
$p\le q\le30$, the formula of Theorem~\ref{thm:main}(b) and the signs in
Theorem~\ref{thm:main}(c) were compared with the coefficients of
$h_{p,q}(1+u)$ computed from Corollary~\ref{cor:numerator}; this checks the
derivation of Theorem~\ref{thm:main}(b) and~(c) from
Corollary~\ref{cor:numerator}.
\end{remark}

Table~\ref{tab:coefficients} lists the Hilbert coefficients for $n\le8$,
computed from Corollary~\ref{cor:numerator}. The
negative entries occur exactly in the rows with $p$ even, in the positions
$p+1,\dots,n-1$.

\begin{table}[ht]
\caption{The Hilbert coefficients $(e_0,\dots,e_n)$ of $I_{p,q}$ for
$p+q\le8$.}\label{tab:coefficients}
\small
\begin{tabular}{@{}cl@{}}
\toprule
$(p,q)$ & $(e_0,e_1,\dots,e_n)$\\
\midrule
$(1,1)$ & $(4,1,0)$\\
$(1,2)$ & $(8,4,1,0)$\\
$(1,3)$ & $(16,12,6,1,0)$\\
$(1,4)$ & $(32,32,24,8,1,0)$\\
$(1,5)$ & $(64,80,80,40,10,1,0)$\\
$(1,6)$ & $(128,192,240,160,60,12,1,0)$\\
$(1,7)$ & $(256,448,672,560,280,84,14,1,0)$\\
$(2,2)$ & $(16,12,1,-2,0)$\\
$(2,3)$ & $(32,32,6,-6,-1,0)$\\
$(2,4)$ & $(64,80,24,-24,-17,-4,0)$\\
$(2,5)$ & $(128,192,80,-80,-90,-32,-3,0)$\\
$(2,6)$ & $(256,448,240,-240,-380,-204,-51,-6,0)$\\
$(3,3)$ & $(64,80,24,1,14,6,0)$\\
$(3,4)$ & $(128,192,80,8,33,20,2,0)$\\
$(3,5)$ & $(256,448,240,40,130,162,78,15,0)$\\
$(4,4)$ & $(256,448,240,40,1,-84,-82,-20,0)$\\
\bottomrule
\end{tabular}
\end{table}

\section*{Use of AI-assisted tools}
AI-assisted tools were used in the research and preparation of this manuscript.

\begin{thebibliography}{9}

\bibitem{BGV}
A.~Banerjee, K.~Goel, and J.~K.~Verma,
\emph{On the Hilbert--Samuel coefficients of Frobenius powers of an ideal},
preprint (2022), arXiv:2203.10173v1.

\bibitem{DNSF}
H.~V.~Do, H.~D.~Nguyen, and S.~A.~Seyed Fakhari,
\emph{On the ordinary and symbolic powers of fiber products},
J. Algebra \textbf{678} (2025), 570--600, \href{https://doi.org/10.1016/j.jalgebra.2025.04.029}{doi:10.1016/j.jalgebra.2025.04.029}.
Numbered references are to arXiv:2409.17551v1.

\bibitem{Fr}
R.~Fr\"oberg,
\emph{Hilbert coefficients of quadratic algebras},
J. Exp. Math. \textbf{1} (2025), no.~2, 315--321, \href{https://doi.org/10.56994/JXM.001.002.007}{doi:10.56994/JXM.001.002.007}.

\bibitem{MN}
A.~Mafi and D.~Naderi,
\emph{Results on the Hilbert coefficients and reduction numbers},
Proc. Indian Acad. Sci. Math. Sci. \textbf{129} (2019), Paper No.~60,
\href{https://doi.org/10.1007/s12044-019-0510-z}{doi:10.1007/s12044-019-0510-z}. Numbered references are to
arXiv:1707.09843v3.

\bibitem{Nar}
M.~Narita,
\emph{A note on the coefficients of Hilbert characteristic functions in
semi-regular local rings},
Proc. Cambridge Philos. Soc. \textbf{59} (1963), 269--275.

\bibitem{NV}
H.~D.~Nguyen and T.~Vu,
\emph{Homological invariants of powers of fiber products},
Acta Math. Vietnam. \textbf{44} (2019), 617--638. Numbered references are to
the corrected version arXiv:1803.04016v3.

\bibitem{Nor}
D.~G.~Northcott,
\emph{A note on the coefficients of the abstract Hilbert function},
J. London Math. Soc. \textbf{35} (1960), 209--214.

\bibitem{Sam}
P.~Samuel,
\emph{La notion de multiplicit\'e en alg\`ebre et en g\'eom\'etrie alg\'ebrique},
J. Math. Pures Appl. (9) \textbf{30} (1951), 159--274.

\end{thebibliography}

\end{document}
